\section{Implementaci\'on y herramientas}

La implementaci\'on se desarroll\'o en Python 3.11 y qued\'o organizada en tres
notebooks finales, ejecutados en orden:

\begin{enumerate}
    \item \texttt{markowitz\_vs\_benchmark\_rebalanceo\_pandemia.ipynb}:
    optimiza carteras Markowitz, calcula el benchmark top-20, aplica rebalanceo
    semestral y compara escenarios con y sin pandemia.

    \item \texttt{black\_litterman.ipynb}: construye portafolios
    Black-Litterman usando las mismas m\'etricas y perfiles, con views de
    momentum, desempleo, momentum top-20 a 6 meses y momentum top-20/bottom-20
    a 1 a\~no.

    \item \texttt{p4\_montecarlo\_recomendador\_portafolios.ipynb}: usa los
    KPIs anteriores como insumo y simula 5\,000 trayectorias para las 180
    combinaciones evaluadas.
\end{enumerate}

Las principales herramientas fueron \textbf{Pandas} para manejo de datos y
exportaci\'on a CSV, \textbf{NumPy/SciPy} para c\'alculo num\'erico y
simulaci\'on, \textbf{CVXPY} para resolver los problemas de optimizaci\'on,
\textbf{Matplotlib} para gr\'aficos, y \textbf{Hashlib} para semillas
deterministas en Monte Carlo.

La trazabilidad se asegura mediante carpetas \texttt{outputs/} por notebook.
Cada etapa exporta CSV con pesos, KPIs, validaciones, rankings o simulaciones,
y esos archivos son los que alimentan las tablas y figuras del informe. Adem\'as,
los notebooks incluyen checklists finales que verifican archivos esperados,
conteos por escenario/t\'ecnica/perfil y ausencia de \texttt{NaN} en columnas
cr\'iticas. No se usan scripts \texttt{.py} externos en la entrega final; el
c\'odigo queda estructurado dentro de los notebooks.
